% ============================================================================
% Fragmento para el TP2 (OE2, Manual y procedimiento funcional de la SDR).
% Viene de los anexos B y C del RE-TP-001 v1.1 (TP1), con el peine de 24
% lineas del EXP-25 y EXP-26. Se integra al main.tex del TP2 cuando se arme
% y no compila solo. Hay que actualizarlo al peine cada 10 MHz.
% ============================================================================

\section{Manejo del software}\label{anexo:sw}

La ventana del flujo muestra en todo momento la etapa, la línea que está saliendo, su portadora y el número de paso dentro de la pasada.

\begin{itemize}
  \item La casilla PAUSA, tildada, deja la salida en silencio. El programa arranca así.
  \item La ganancia de TX tiene tres valores, 45, 55 y 65~dB, que dejan las líneas calibradas en $-20$, $-10$ y 0~dBm. La calibración vale para los tres.
  \item El nivel puede ser Calibrado o Calibrado + 3~dB. La segunda opción sube todas las líneas, y en la medición por conducción dieron +2,3~dBm de media porque la placa ya comprime un poco.
  \item Importar calibración previa carga en los 24 contadores los valores del archivo de calibración.
  \item Guardar calibración escribe en el archivo los valores que están en los contadores.
  \item Generar calibración arma la calibración a partir de una captura del analizador (sección~\ref{anexo:cal}).
  \item Hay un contador por línea, que guarda el nivel leído con el peine sin calibrar menos 20~dB. Para bajar una línea $x$~dB se le suma $x$ a su contador, y para subirla se le resta.
\end{itemize}

El 3.\textsuperscript{er} LED de la SDR, contando de izquierda a derecha, debe quedar encendido fijo mientras el flujo corre. Si titila, a la placa le están faltando muestras por el USB y la medición no es válida.

\section{Calibración de la fuente por conducción}\label{anexo:cal}

Se hace con la SDR conectada por cable al analizador portátil, con atenuación de 30~dB, de 45 a 1005~MHz, RBW de 300~kHz y MAX HOLD.

\begin{enumerate}
  \item Abrir el flujo y ejecutarlo sin importar ninguna calibración. Los 24 contadores deben estar en cero.
  \item Elegir ganancia 65~dB y destildar PAUSA.
  \item Reiniciar MAX HOLD, esperar una pasada completa y guardar la captura en formato .csv.
  \item Reiniciar MAX HOLD y guardar una segunda captura de otra pasada. Con dos capturas el programa descarta los picos momentáneos.
  \item Apretar Generar calibración y elegir las dos capturas a la vez. El programa busca el pico de cada línea, carga los contadores y guarda el archivo de calibración.
  \item Reiniciar MAX HOLD, esperar una pasada y verificar que las 24 líneas queden en 0~dBm~$\pm$~1~dB. Si alguna no lo cumple, corregir su contador a mano y apretar Guardar calibración.
\end{enumerate}

La calibración vale para el cable y el analizador con que se hizo. Lleva unos 20~minutos.
